\documentclass[10pt]{article}
\usepackage[T1]{fontenc}
\usepackage[utf8]{inputenc}
\usepackage[french]{babel}
\usepackage{amsmath,amssymb,mathrsfs,mathpazo}
\usepackage{geometry}
\geometry{a4paper,margin=2cm}
\usepackage{graphicx,array,tabularx,tabularray,colortbl,multirow,multicol,booktabs,enumitem,pifont,ulem,fancybox,diagbox,cancel}
\usepackage{pstricks,pstricks-add,pst-plot,pst-tree,pst-node,pst-text,pst-func,pst-math,pst-bspline,pst-3dplot,pst-eucl,pst-solides3d}
\let\pscalculate\undefined
\usepackage{pst-bezier}
\usepackage{tikz,pgfplots,tkz-tab}
\usetikzlibrary{arrows,arrows.meta,patterns,calc,positioning,trees,decorations.markings,shapes.geometric,shapes.misc,shapes.arrows,intersections,decorations.pathmorphing,decorations.pathreplacing,backgrounds,shadows,babel}
\pgfplotsset{compat=1.18}
\usepgfplotslibrary{fillbetween}
\usepackage[np]{numprint}
\usepackage{listings,makecell,pas-tableur,twoopt}
\let\barre\undefined
\usepackage[pstricks]{tablvar}
\definecolor{darkgreen}{RGB}{0,100,0}
\providecommand{\red}{\color{red}}
\providecommand{\blue}{\color{blue}}
\providecommand{\green}{\color{green}}
\providecommand{\black}{\color{black}}
\providecommand{\white}{\color{white}}
\providecommand{\gray}{\color{gray}}
\providecommand{\cyan}{\color{cyan}}
\providecommand{\magenta}{\color{magenta}}
\providecommand{\yellow}{\color{yellow}}
\providecommand{\violet}{\color{violet}}
\providecommand{\orange}{\color{orange}}
\providecommand{\euro}{\texteuro}
\providecommand{\EUR}{\texteuro}
\providecommand{\ieme}{\textsuperscript{e}}
\providecommand{\degres}{\ensuremath{^\circ}}
\providecommand{\degree}{\ensuremath{^\circ}}
\providecommand{\vv}[1]{\overrightarrow{#1}}
\providecommand{\Coord}[1]{#1}
\providecommand{\decosix}{}
\providecommand{\starredbullet}{\ensuremath{\bullet}}
\pagestyle{empty}
\setlength{\parindent}{0pt}
\usepackage{adjustbox}\DeclareUnicodeCharacter{2236}{:}\DeclareUnicodeCharacter{2190}{\ensuremath{\leftarrow}}
\providecommand{\R}{\mathbb{R}}
\providecommand{\N}{\mathbb{N}}
\providecommand{\Z}{\mathbb{Z}}
\providecommand{\Q}{\mathbb{Q}}
\providecommand{\D}{\mathbb{D}}
\providecommand{\C}{\mathbb{C}}
\providecommand{\e}{{\mathrm{e}}}
\providecommand{\vect}[1]{\overrightarrow{#1}}
\providecommand{\vectt}[1]{\overrightarrow{\mathrm{#1}}}
\providecommand{\vvt}[1]{\overrightarrow{\mathrm{#1}}}
\providecommand{\np}[2][]{{#2~#1}}
\providecommand{\vv}[1]{{\overrightarrow{#1}}}
\renewcommand{\d}{{\mathrm{d}}}
\providecommand{\degres}{{^\circ}}
\providecommand{\degre}{{^\circ}}
\providecommand{\blue}{\color{blue}}
\providecommand{\red}{\color{red}}
\providecommand{\enskip}{\quad}
\providecommand{\llbracket}{\lbrack\!\lbrack}
\providecommand{\rrbracket}{\rbrack\!\rbrack}
\providecommand{\euro}{{\text{€}}}
\providecommand{\no}{{\text{n°}}}
\providecommand{\interval}[2]{{[#1;#2]}}
\providecommand{\intervFO}[2]{{[#1;#2[}}
\providecommand{\dotfill}{\ldots}

\providecommand{\ds}{\displaystyle}
\providecommand{\textding}[1]{#1}
\providecommand{\Oij}{\(\mathrm{O};\vec{\imath},\vec{\jmath}\)}
\providecommand{\Oijk}{\(\mathrm{O};\vec{\imath},\vec{\jmath},\vec{k}\)}
\providecommand{\Ouv}{\(\mathrm{O};\vec{u},\vec{v}\)}
\providecommand{\intervFF}[2]{\ensuremath{[#1;#2]}}
\renewcommand{\euro}{\ifmmode\text{\texteuro}\else\texteuro\fi}
\DeclareUnicodeCharacter{00B0}{\ifmmode^{\circ}\else\textdegree\fi}
\DeclareUnicodeCharacter{20AC}{\ifmmode\text{\texteuro}\else\texteuro\fi}
% Énoncés : Éducation nationale. Transcription et corrigés : APMEP, Denis Vergès et contributeurs. Découpage et classement : CoopMaths, Sébastien Lozano, Éric Elter et contributeurs. Adaptation de la présentation : Maths Explorer.
% CC BY-SA : https://coopmaths.fr/www/about/
% https://coopmaths.fr/alea/static/bac/2021/tex/bac_2021_09_metropole_10.tex
\begin{document}
\noindent\textit{APMEP / CoopMaths. CC BY-SA. Présentation adaptée par Maths Explorer.}\par\medskip
\medskip

\begin{tabular}{|l|}\hline
Principaux domaines abordés :\\
Fonction logarithme, limites, dérivation.\\ \hline
\end{tabular}

\medskip

\textbf{Partie 1}

\medskip

Le graphique ci-dessous donne la représentation graphique dans un repère orthonormé de la fonction $f$
définie sur l'intervalle $]0~;~+\infty[$ par:

\[f(x) = \dfrac{2\ln (x) - 1}{x}.\]

\begin{center}
\psset{unit=1.75cm}
\begin{pspicture*}(-0.6,-2.2)(5,1)
\psaxes[linewidth=1.25pt,labelFontSize=\scriptstyle]{->}(0,0)(0,-2.1)(5,1)
\psplot[plotpoints=2000,linewidth=1.25pt,linecolor=red]{0.1}{5}{x ln 2 mul 1 sub x div}
\end{pspicture*}
\end{center}

\medskip

\begin{enumerate}
\item Déterminer par le calcul l'unique solution $\alpha$ de l'équation $f(x) = 0$.

On donnera la valeur exacte de $\alpha$ ainsi que la valeur arrondie au centième.
\item Préciser, par lecture graphique, le signe de $f(x)$ lorsque $x$ varie dans l'intervalle $]0~;~+\infty[$.
\end{enumerate}

\bigskip

\textbf{Partie II}

\medskip
On considère la fonction $g$ définie sur l'intervalle $]0~;~+\infty[$ par:

\[g(x) = [\ln (x)]^2 - \ln (x).\]

\begin{enumerate}
\item 
	\begin{enumerate}
		\item Déterminer la limite de la fonction $g$ en $0$.
		\item Déterminer la limite de la fonction $g$ en $+ \infty$.
	\end{enumerate}
\item On note $g'$ la fonction dérivée de la fonction $g$ sur l'intervalle $]0~;~+\infty[$.

Démontrer que, pour tout nombre réel $x$ de $]0~;~+\infty[$, on a : $g'(x) = f(x)$, où $f$ désigne la fonction définie dans la partie I.
\item Dresser le tableau de variations de la fonction $g$ sur l'intervalle $]0~;~+\infty[$.

On fera figurer dans ce tableau les limites de la fonction $g$ en $0$ et en $+\infty$, ainsi que la valeur du minimum de $g$ sur $]0~;~+\infty[$.
\item Démontrer que, pour tout nombre réel $m > - 0,25$, l'équation $g(x) = m$ admet exactement deux solutions.
\item Déterminer par le calcul les deux solutions de l'équation $g(x) = 0$.
\end{enumerate}

\end{document}
